\documentclass[12pt,a4paper]{article}

\usepackage[utf8]{inputenc}
\usepackage[T1]{fontenc}
\usepackage[english]{babel}
\usepackage{amsmath, amssymb, amsthm}
\usepackage{bm}
\usepackage{graphicx}
\usepackage{booktabs}
\usepackage{hyperref}
\usepackage{geometry}
\usepackage{microtype}
\usepackage{xcolor}
\usepackage{algorithm}
\usepackage{algpseudocode}
\usepackage{caption}
\usepackage{subcaption}
\usepackage{enumitem}

\geometry{margin=2.5cm}

\hypersetup{
    colorlinks=true,
    linkcolor=blue!60!black,
    citecolor=green!50!black,
    urlcolor=blue!70!black,
}

% ── math shorthands ──────────────────────────────────────────────────────────
\newcommand{\R}{\mathbb{R}}
\newcommand{\kap}{\kappa}
\newcommand{\Jprod}[2]{\langle #1,\, #2 \rangle_J}
\newcommand{\phipos}{\varphi_+}
\newcommand{\phineg}{\varphi_-}
\newcommand{\dpoly}{d_{\mathrm{poly}}}

\theoremstyle{definition}
\newtheorem{definition}{Definition}
\newtheorem{remark}{Remark}

% ─────────────────────────────────────────────────────────────────────────────
\title{%
  \textbf{Semantic Segmentation in Pontryagin Spaces}\\
  \large Random Feature Embeddings with Indefinite Metric for UAV Crop Imagery
}
\author{Carlos Jiménez}
\date{\today}

\begin{document}
\maketitle
\tableofcontents
\newpage

% =============================================================================
\section{Introduction}
% =============================================================================

Deep neural network classifiers typically embed intermediate representations
in Euclidean space and apply a linear decision rule.  A growing body of work
has shown that certain data distributions — in particular those with latent
hierarchical or hyperbolic structure — are better captured by non-Euclidean
geometries.  Hyperbolic spaces (Poincaré ball, hyperboloid) have been
explored for this purpose~\cite{ganea2018,atigh2022}, but they impose a
\emph{constant negative} curvature that may not suit all problems.

This work proposes an alternative: embedding pixel features into a
\textbf{Pontryagin space} $\Pi_\kap$, a flat pseudo-Riemannian manifold
equipped with an \emph{indefinite} inner product of signature $(\cdot, \kap)$.
Unlike the Poincaré ball, a Pontryagin space is both flat and able to
represent mixed-curvature structure through the coexistence of positive-
and negative-definite subspaces.  The embedding is constructed from two
families of random features:

\begin{itemize}[noitemsep]
  \item \textbf{Random Fourier Features (RFF)}~\cite{rahimi2007} approximate
        a positive-definite RBF kernel $K_+$.
  \item \textbf{Spherical Random Features (SRF)}~\cite{pennington2015}
        approximate a polynomial kernel $K_-$ of degree $\dpoly$ on the unit
        sphere.
\end{itemize}

The resulting feature vector lives in $\Pi_\kap = \R^{p+\kap}$ with metric
$J = \mathrm{diag}(+1^p, -1^\kap)$, and classification is performed by
\emph{J-hyperplane} multinomial logistic regression.

The experimental testbed is binary semantic segmentation of sugarcane crops
from UAV imagery, where all three heads (\textsc{Euclidean},
\textsc{Hyperbolic}, \textsc{Pontryagin}) share an identical deep FCN backbone
and training protocol, isolating the effect of the embedding geometry.
Interpretability of each model is assessed using \textbf{SegScoreCAM},
an adaptation of ScoreCAM for pixel-wise segmentation outputs.

% =============================================================================
\section{Dataset}
% =============================================================================

\subsection{UAV Sugarcane Segmentation}

The dataset consists of aerial images acquired by an unmanned aerial vehicle
(UAV) over agricultural fields containing sugarcane, banana trees, pepper,
spinach, and other crops.  The annotation format is \textbf{YOLO polygon
segmentation}: each label file contains class index and normalised polygon
vertices per instance.

\paragraph{Preprocessing.}
All images are resized to $256 \times 256$ pixels using bilinear
interpolation.  Polygon masks are rasterised at the same resolution.
To focus on a single, well-defined binary problem, the dataset is filtered
to \textbf{binary sugarcane segmentation}: class~0 = background, class~1 =
sugarcane; images without any sugarcane polygon are discarded.
Pixel intensities are normalised with ImageNet statistics
($\mu = (0.485, 0.456, 0.406)$, $\sigma = (0.229, 0.224, 0.225)$).

\paragraph{Class imbalance.}
Background pixels are far more frequent than sugarcane pixels.  A weighted
cross-entropy loss compensates for this imbalance with weights
$w_{\text{bg}} = 0.20$ and $w_{\text{sc}} = 1.00$.

\paragraph{Splits and augmentation.}
The dataset is split into train / validation / test partitions provided by
the source repository, containing 81 / 30 / 20 images respectively (before
binary filtering).  After discarding images without sugarcane, the effective
splits used in training are 49 train / 23 validation / 13 test images.
During training, random horizontal and vertical flips (each with probability
0.5) are applied.

% =============================================================================
\section{Backbone Architectures}
% =============================================================================

Two backbone architectures are evaluated: a lightweight deep FCN for rapid
ablation and a full U-Net for competitive benchmarking.  All segmentation
heads (Euclidean, Hyperbolic, Pontryagin) are compatible with both backbones;
the U-Net experiments additionally include a Vanilla baseline that trains
without class-weighted loss.

\subsection{DeepFCN}
\label{sec:deepfcn}

All FCN segmentation heads share a \textbf{deep fully-convolutional network
(DeepFCN)} backbone with configurable depth $D \in \{3,4,5\}$.  Each stage
consists of two $3\times3$ strided convolutions (stride 2) with batch
normalisation and ReLU activation.  Channel widths follow the schedule
$3 \to 32 \to 64 \to 128 \to 256 \to 256$ (the last stage is repeated when
$D = 5$).  A final $1 \times 1$ projection maps the last stage to a fixed
$C_{\text{out}} = 64$ channel dimension.

Given an input $\mathbf{x} \in \R^{3 \times H \times W}$, the backbone
produces a feature map $\mathbf{f} \in \R^{64 \times H' \times W'}$ where
$H' = H / 2^D$.  This map is then processed by the segmentation head, and
the resulting low-resolution logits are bilinearly upsampled to the original
resolution before computing the loss.

\subsection{U-Net Backbone}
\label{sec:unet}

The U-Net backbone (Ronneberger et al., 2015) is an encoder-decoder
architecture that preserves full spatial resolution through skip connections,
avoiding the upsampling artefacts common in FCN-style designs.

\paragraph{Encoder.}
$D$ double-conv blocks (two $3\times3$ convolutions, each followed by batch
normalisation and ReLU), each followed by $2\times2$ max-pooling.  The
channel schedule is $C_k = C_{\text{base}} \cdot \min(2^k, 8)$ for stage
$k = 0,\ldots,D-1$, where $C_{\text{base}} = 64$ by default.  For $D = 4$
this gives $64 \to 128 \to 256 \to 512$.

\paragraph{Bottleneck.}
A single double-conv block at the lowest resolution:
$C_{D-1} \to C_D = C_{\text{base}} \cdot \min(2^D, 8)$
(e.g.\ $512 \to 1024$ for $D = 4$).

\paragraph{Decoder.}
$D$ upsampling stages, each performing bilinear upsampling followed by
skip-connection concatenation with the matching encoder output, and then a
double-conv block.  The output is projected to $C_{\text{out}} = 64$
channels by a $1 \times 1$ convolution.

The backbone produces $\mathbf{f} \in \R^{64 \times H \times W}$ — at the
\emph{same} spatial resolution as the input, so no further upsampling is
needed.  The depth $D \in \{3,4,5\}$ is a hyperparameter optimised in the
HPO phase.

% =============================================================================
\section{Segmentation Heads}
% =============================================================================

The same four heads are used with both backbones.  The Vanilla head is a
pure CE baseline without class weighting; it is included in the U-Net
experiments to provide an upper bound untouched by any geometric inductive
bias.

\subsection{Vanilla (U-Net baseline only)}\label{sec:vanilla}

The Vanilla head applies a $1\times1$ convolution to the backbone features
and minimises the standard (unweighted) cross-entropy loss:
\[
  \mathcal{L}_{\text{vanilla}}
  = -\frac{1}{N}\sum_{n=1}^{N}\log\frac{e^{(\mathbf{W}_{\text{cls}}\mathbf{f}_n)_{y_n}}}
    {\sum_c e^{(\mathbf{W}_{\text{cls}}\mathbf{f}_n)_c}}.
\]
No class-weight correction is applied.  This model is the strongest
``geometry-free'' reference: any improvement from the geometric heads must
overcome a baseline that is otherwise identical.

\subsection{Euclidean (Baseline)}

The Euclidean head applies a $1 \times 1$ convolution directly to the
backbone features:
\[
  \hat{\mathbf{y}} = \mathbf{W}_{\text{cls}}\, \mathbf{f} + \mathbf{b},
  \quad \mathbf{W}_{\text{cls}} \in \R^{C \times 64}.
\]
This is the standard FCN classifier; no geometric embedding is applied.

\subsection{Hyperbolic (Ganea / Atigh Baseline)}

The hyperbolic head maps backbone features into the Poincaré ball
$\mathbb{B}^d_c$ of curvature $c$ using a Möbius transformation, then
applies the hyperbolic multinomial logistic regression (MLR) of
Ganea et al.~\cite{ganea2018} extended to the pixel level by Atigh et
al.~\cite{atigh2022}.  The decision boundary for class $y$ is a
\emph{geodesic hyperplane} in $\mathbb{B}^d_c$.

The curvature $c$ is a hyperparameter optimised during search.

\subsection{Pontryagin (This Work)}

\paragraph{Motivation.}
The Poincaré ball has constant negative curvature, which biases the space
towards tree-like structures.  A Pontryagin space $\Pi_\kap$ is flat but
has an \emph{indefinite} metric; it can simultaneously encode positive-
and negative-curvature regions through the balance of its positive and
negative subspace dimensions.

\begin{definition}[Pontryagin space $\Pi_\kap$]
  A Pontryagin space of signature $\kap$ is a real vector space $\R^{p+\kap}$
  equipped with the indefinite inner product
  \[
    \Jprod{\mathbf{u}}{\mathbf{v}}
    = \sum_{i=1}^{p} u_i v_i - \sum_{i=p+1}^{p+\kap} u_i v_i
    = \mathbf{u}^\top J\, \mathbf{v},
    \quad J = \mathrm{diag}(\underbrace{+1,\ldots,+1}_{p},\,
                             \underbrace{-1,\ldots,-1}_{\kap}).
  \]
\end{definition}

\paragraph{Random feature construction.}
Given a backbone feature vector $\mathbf{f} \in \R^{64}$ at each pixel, the
Pontryagin embedding is built as follows.

\subparagraph{Positive component — RFF.}
Random Fourier Features~\cite{rahimi2007} approximate the RBF kernel
\[
  K_+(\mathbf{x}, \mathbf{y}) = \exp\!\left(-\tfrac{\|\mathbf{x}-\mathbf{y}\|^2}{2\sigma^2}\right).
\]
Frequencies $\omega_i \sim \mathcal{N}(0, \sigma^{-2} I)$ and phases
$b_i \sim \mathrm{Uniform}[0, 2\pi]$ are sampled once; the feature map is
\[
  \phipos(\mathbf{x}) = \frac{1}{\sqrt{n_{\text{rff}}}}
  \Bigl[\cos(\omega_1^\top \mathbf{x} + b_1),\, \sin(\omega_1^\top \mathbf{x} + b_1),\,
        \ldots\Bigr]
  \in \R^{2 n_{\text{rff}}}.
\]
The bandwidth $\sigma$ is stored as a learnable parameter $\log\sigma$, which
is updated by gradient descent and therefore \emph{removed from the HPO
space}.

\subparagraph{Negative component — SRF.}
Spherical Random Features~\cite{pennington2015} approximate the polynomial
kernel
\[
  K_-(\mathbf{x}, \mathbf{y}) = (\mathbf{x}^\top \mathbf{y})^{\dpoly},
\]
where $\dpoly$ is the polynomial degree (\emph{independent of the signature
$\kap$}).  The $j$-th feature is
\[
  \phineg(\mathbf{x})_j = \frac{1}{\sqrt{\kap}}
  \prod_{\ell=1}^{\dpoly} \bigl(\omega_{j\ell}^\top \mathbf{x}\bigr),
  \quad
  \omega_{j\ell} \sim \mathrm{Uniform}(S^{d-1}),
\]
producing $\phineg(\mathbf{x}) \in \R^\kap$.

\subparagraph{Concatenation.}
The Pontryagin embedding is
\[
  \Phi(\mathbf{x})
  = \bigl[\phipos(\mathbf{x})\,\|\,\phineg(\mathbf{x})\bigr]
  \in \R^{p + \kap},
  \quad p = 2\,n_{\text{rff}}.
\]
The metric matrix $J = \mathrm{diag}(+1^p, -1^\kap)$ is stored as a fixed
buffer.

\begin{remark}
  The signature $\kap$ and the polynomial degree $\dpoly$ are
  \textbf{decoupled hyperparameters}: $\kap$ sets how many negative
  directions exist in $J$ (and therefore how many SRF components are
  produced), while $\dpoly$ controls the nonlinearity of the kernel $K_-$.
  Both are optimised independently in the HPO search.
\end{remark}

\paragraph{J-hyperplane classification.}
For each class $y$, the logit is the indefinite inner product with a learnable
weight vector $\mathbf{W}_y \in \R^{p+\kap}$:
\[
  \zeta_y(\mathbf{z}) = \Jprod{\mathbf{W}_y}{\mathbf{z}} + b_y
  = \sum_{d=1}^{p+\kap} J_d\, W_{y,d}\, z_d + b_y.
\]

% =============================================================================
\section{Loss Function}
% =============================================================================

The Pontryagin head minimises a composite loss:
\[
  \mathcal{L} = \mathcal{L}_{\text{CE}}
              + \lambda_{\text{topo}}\, \mathcal{L}_{\text{topo}}
              + \lambda_{\text{bal}}\,  \mathcal{L}_{\text{balance}}.
\]

\begin{remark}[Cone penalty disabled]
An earlier version of the loss included a cone penalty on the classifier
weight vectors:
\[
  \mathcal{L}_{\text{cone-W}}
  = \frac{1}{C}\sum_{y=1}^C
    \mathrm{ReLU}\!\left(\varepsilon - \bigl|\mathbf{W}_y^\top J\, \mathbf{W}_y\bigr|\right).
\]
Empirically, this term suppresses the negative subspace by keeping all
weight vectors strictly timelike, which defeats the purpose of the
indefinite metric.  It has therefore been removed from the active loss
($\lambda_W = 0$); the balance penalty described below serves instead to
maintain gradient flow in the negative dimensions.
\end{remark}

\paragraph{Cross-entropy.}
\[
  \mathcal{L}_{\text{CE}} =
  -\frac{1}{N}\sum_{n=1}^N \log \frac{e^{\zeta_{y_n}(\mathbf{z}_n)}}
                                     {\sum_c e^{\zeta_c(\mathbf{z}_n)}},
\]
weighted by class frequencies ($w_{\text{bg}}=0.20$, $w_{\text{sc}}=1.00$).

\paragraph{Topographic penalty $\mathcal{L}_{\text{topo}}$.}
The \texttt{IndefiniteTopographicPenalty} combines three terms on the
embedding $\mathbf{z}$:
\begin{itemize}[noitemsep]
  \item \emph{Light-cone (LC):} penalises embeddings near the null cone
        $\{\mathbf{z} : \Jprod{\mathbf{z}}{\mathbf{z}} = 0\}$.
  \item \emph{Causal consistency (CC):} encourages same-class embeddings to
        share the same causal character (timelike or spacelike).
  \item \emph{Signature balance (SB):} keeps the norms of the positive and
        negative subcomponents similar.
\end{itemize}
Weights $(\lambda_{\text{lc}}, \lambda_{\text{cc}}, \lambda_{\text{sb}}) =
(1.0,\, 0.5,\, 0.5)$ are held fixed; $\lambda_{\text{topo}}$ scales the
whole term and is searched during HPO.

\paragraph{Balance penalty $\mathcal{L}_{\text{balance}}$.}
\[
  \mathcal{L}_{\text{balance}}
  = \bigl(\mathrm{rms}(\mathbf{W}_+) - \mathrm{rms}(\mathbf{W}_-)\bigr)^2,
\]
where $\mathbf{W}_+$ and $\mathbf{W}_-$ are the positive- and
negative-subspace columns of the classifier weight matrix, and
$\mathrm{rms}(\cdot)$ denotes the root mean square.  This penalty encourages
the two subspaces to carry equal energy in the weight vectors, ensuring that
neither the positive (RFF) nor the negative (SRF) subspace dominates
gradient flow during training.

% =============================================================================
\section{Training Protocol}
% =============================================================================

Both experiment tracks (FCN and U-Net) share the same protocol.

\paragraph{Optimiser.}  Adam with learning rate $\eta = 10^{-3}$ and weight
decay $10^{-4}$ (default values, overridden by the HPO best parameters when
available).

\paragraph{Scheduler.}  Cosine annealing from $\eta$ to
$\eta_{\min} = 10^{-5}$ over $E_{\max} = 60$ epochs.

\paragraph{Early stopping.}  Training halts if the validation \emph{macro
mIoU} (averaged over non-empty classes) does not improve for 15 consecutive
epochs.  The checkpoint with the highest validation macro mIoU is restored
for test evaluation.

\paragraph{Batch size.}  8 images per GPU step at $256\times256$ resolution.

\paragraph{Hardware.}  Single NVIDIA A100-SXM4 (80 GB) GPU on Google Colab.

% =============================================================================
\section{Hyperparameter Optimisation}
% =============================================================================

Bayesian HPO is carried out with \texttt{scikit-optimize}'s
\texttt{forest\_minimize} (Extra Trees surrogate, 30 trials, 10 random
starts).  Each trial trains for 15 epochs (FCN) or 12 epochs (U-Net).
The search spaces differ between the two backbones only in the
architecture-depth parameter; all other parameters are shared.

\subsection{FCN experiments}
\label{sec:hpo_fcn}

\begin{center}
\begin{tabular}{llll}
\toprule
Parameter & Range & Prior & Applies to \\
\midrule
\texttt{lr}              & $[10^{-5},\; 10^{-2}]$ & log-uniform & all \\
\texttt{weight\_decay}   & $[10^{-6},\; 10^{-3}]$ & log-uniform & all \\
\texttt{backbone\_depth} & $\{3, 4, 5\}$           & uniform     & all \\
\texttt{hyperbolic\_c}   & $[0.1,\; 2.0]$          & log-uniform & hyperbolic \\
$\kap$                   & $\{1,\ldots,64\}$       & log-uniform & pontryagin \\
$\dpoly$                 & $\{1,\ldots,10\}$       & uniform     & pontryagin \\
\texttt{rff\_multiplier} & $\{1,2,3,4\}$           & uniform     & pontryagin \\
$\lambda_{\text{topo}}$  & $[10^{-6},\; 0.2]$      & log-uniform & pontryagin \\
$\lambda_{\text{bal}}$   & $[10^{-3},\; 1.0]$      & log-uniform & pontryagin \\
\bottomrule
\end{tabular}
\end{center}

\subsection{U-Net experiments}
\label{sec:hpo_unet}

The U-Net HPO space is identical to the FCN one with a single change:
\texttt{backbone\_depth} is replaced by \texttt{unet\_depth}, which
controls the number of encoder/decoder stages of the U-Net backbone
(Section~\ref{sec:unet}).

\begin{center}
\begin{tabular}{llll}
\toprule
Parameter & Range & Prior & Applies to \\
\midrule
\texttt{lr}              & $[10^{-5},\; 10^{-2}]$ & log-uniform & all \\
\texttt{weight\_decay}   & $[10^{-6},\; 10^{-3}]$ & log-uniform & all \\
\texttt{unet\_depth}     & $\{3, 4, 5\}$           & uniform     & all \\
\texttt{hyperbolic\_c}   & $[0.1,\; 2.0]$          & log-uniform & hyperbolic \\
$\kap$                   & $\{1,\ldots,64\}$       & log-uniform & pontryagin \\
$\dpoly$                 & $\{1,\ldots,10\}$       & uniform     & pontryagin \\
\texttt{rff\_multiplier} & $\{1,2,3,4\}$           & uniform     & pontryagin \\
$\lambda_{\text{topo}}$  & $[10^{-6},\; 0.2]$      & log-uniform & pontryagin \\
$\lambda_{\text{bal}}$   & $[10^{-3},\; 1.0]$      & log-uniform & pontryagin \\
\bottomrule
\end{tabular}
\end{center}

\paragraph{Design rationale for $\kap$.}
The signature $\kap$ determines how many negative dimensions exist in $J$ and
simultaneously the number of SRF output components.  A log-uniform prior over
$\{1,\ldots,64\}$ ensures that small values ($\kap \leq 8$, where the
indefinite component is modest) receive the same sampling probability as large
values ($\kap \geq 32$, where the negative subspace is dominant).

\paragraph{Decoupling $\kap$ and $\dpoly$.}
In earlier implementations, $\kap$ controlled both the number of SRF
components and the polynomial degree of $K_-$, conflating three distinct
design choices.  The current implementation separates them:
$\kap = n_{\text{srf}}$ sets the signature, while $\dpoly$ sets the degree
of $K_-(\mathbf{x},\mathbf{y}) = (\mathbf{x}^\top\mathbf{y})^{\dpoly}$,
searched independently.

\paragraph{Gradient-tuned parameters.}
The RFF bandwidth $\sigma$ is \emph{not} part of the HPO space.  It is
reparametrised as $\sigma = \exp(\log\sigma)$ and updated by gradient
descent, with unit-variance frequencies $\omega_i \sim \mathcal{N}(0,I)$
rescaled by $1/\sigma$ at each forward pass.

% =============================================================================
\section{Evaluation Metrics}
% =============================================================================

All metrics are computed on the held-out test split.

\paragraph{Per-class metrics.}
For each class $c$: true positives ($\mathrm{TP}_c$), false positives
($\mathrm{FP}_c$), false negatives ($\mathrm{FN}_c$), and true negatives
($\mathrm{TN}_c$); then
\[
  \mathrm{IoU}_c = \frac{\mathrm{TP}_c}{\mathrm{TP}_c + \mathrm{FP}_c + \mathrm{FN}_c},
  \qquad
  \mathrm{Dice}_c = \frac{2\,\mathrm{TP}_c}{2\,\mathrm{TP}_c + \mathrm{FP}_c + \mathrm{FN}_c}.
\]
Precision, recall, and specificity (TNR) are also reported.

\paragraph{Global metrics.}
\begin{itemize}[noitemsep]
  \item \emph{Macro mIoU / Dice}: average over non-empty classes.
  \item \emph{Micro IoU / Dice}: pooled across all pixels.
  \item \emph{Pixel accuracy}: fraction of correctly labelled pixels.
\end{itemize}

The monitor criterion for early stopping and HPO is the \textbf{validation
macro mIoU}.

% =============================================================================
\section{Interpretability: SegScoreCAM}
% =============================================================================

\subsection{Motivation}

The original ScoreCAM~\cite{scorecam2020} was designed for image
\emph{classification}: the score for each feature-map channel is the scalar
class logit after masking.  Segmentation heads produce spatial output maps
$\hat{P} \in \R^{C \times H \times W}$; extracting a single logit discards
all spatial information.

\subsection{SegScoreCAM}

We replace the scalar score with the \textbf{mean softmax probability} of the
target class over all spatial positions:

\[
  s_k = \frac{1}{H W}
        \sum_{i,j} \mathrm{softmax}\bigl(f_\theta(I \odot \tilde{A}_k)\bigr)_{c,i,j},
\]

where $\tilde{A}_k$ is the $k$-th feature map normalised to $[0,1]$ and
upsampled to input resolution, and $I \odot \tilde{A}_k$ is the masked input.

The class activation map for class $c$ is then:

\[
  \mathrm{CAM}_c = \mathrm{ReLU}\!\left(
    \sum_k \bar{s}_k \cdot A_k
  \right), \quad
  \bar{s}_k = \frac{|s_k|}{\sum_{k'} |s_{k'}|},
\]

followed by min-max normalisation to $[0,1]$.

\subsection{Interpretability Metrics}
\label{sec:cam_metrics}

All metrics are computed on the sugarcane class (class~1), where the
ground-truth mask $M \in \{0,1\}^{H \times W}$ is available.

\paragraph{Adaptive binarisation (Otsu).}
Earlier evaluations binarised the CAM at a fixed threshold of $0.5$.  This
is inappropriate when heatmaps have very different dynamic ranges across
models.  We replace it with \textbf{Otsu's method}: the threshold
$\tau^* \in [0,1]$ that maximises the inter-class variance of the CAM
histogram, computed independently for each image.  Let
$\hat{M} = \mathbf{1}[\mathrm{CAM} \geq \tau^*]$ be the resulting binary map.

\paragraph{GT-alignment metrics.}
\begin{center}
\begin{tabular}{lll}
\toprule
Metric & Formula & Interpretation \\
\midrule
GT IoU
  & $\dfrac{|\hat{M} \cap M|}{|\hat{M} \cup M|}$
  & overlap with GT ↑ \\[6pt]
GT Recall
  & $\dfrac{|\hat{M} \cap M|}{|M|}$
  & coverage of GT pixels ↑ \\[6pt]
GT Precision
  & $\dfrac{|\hat{M} \cap M|}{|\hat{M}|}$
  & fraction of CAM on GT ↑ \\[6pt]
GT Pearson
  & $\rho\!\left(\mathrm{CAM}_{:},\; M_{:}\right)$
  & linear correlation (no threshold) ↑ \\[4pt]
Mean activation
  & $\frac{1}{HW}\sum_{i,j}\mathrm{CAM}_{i,j}$
  & diffuseness ↓ \\
\bottomrule
\end{tabular}
\end{center}

\paragraph{Confidence gain.}
\label{sec:conf_gain}
Beyond spatial overlap with the GT mask, we assess whether the highlighted
region is \emph{causally sufficient} for the model's prediction by measuring
how model confidence changes when only the explanatory region is shown.
Given the explanatory image $I_{\text{exp}} = I \odot \mathrm{CAM}$, we
compute the per-pixel percentage change in softmax confidence for class $c$:

\[
  \Delta_{i,j}^{(c)}
  = \frac{P_{\text{exp}}(i,j,c) - P_{\text{orig}}(i,j,c)}
         {P_{\text{orig}}(i,j,c) + \epsilon} \times 100,
\]

where $P_{\text{orig}}$ and $P_{\text{exp}}$ are the softmax outputs on the
original and masked images respectively.  We then restrict to pixels where
the GT label equals $c$ and report the \emph{mean} and \emph{standard
deviation} of $\Delta^{(c)}$ over those pixels.  A \textbf{negative} mean
indicates that masking the input by the CAM reduces confidence on the
relevant pixels — the model loses information it needs.  A value close to
zero or positive indicates that the highlighted region is self-sufficient for
the prediction.

\subsection{Interpretability Results}

Table~\ref{tab:cam_full} reports all metrics for the sugarcane class.
Values marked with $(\star)$ are sample estimates from the first recorded
test images; full-dataset aggregates require re-running
\texttt{interpretability\_analysis.py} after the HPO re-training with the
expanded $\kap$ search space.

\begin{table}[h]
\centering
\caption{SegScoreCAM interpretability metrics on the sugarcane class
  (sample values from available test images, new metric suite with Otsu
  binarisation).  All GT-alignment metrics are higher-is-better;
  Mean Activation is lower-is-better.  Confidence gain is the mean \%
  change in softmax probability on GT-positive pixels when the input is
  masked by the CAM.}
\label{tab:cam_full}
\resizebox{\textwidth}{!}{%
\begin{tabular}{lcccccc}
\toprule
Model
  & GT IoU $\uparrow$
  & GT Recall $\uparrow$
  & GT Precision $\uparrow$
  & GT Pearson $\uparrow$
  & Mean Act.\ $\downarrow$
  & Conf.\ gain (\%) \\
\midrule
Euclidean
  & $0.658$ & $0.764$ & $0.827$ & $0.836$
  & $0.040$ & $-87.0 \pm 27.0$ \\
Hyperbolic
  & $0.549$ & $0.596$ & $0.874$ & $0.787$
  & $0.029$ & $-40.6 \pm 31.0$ \\
Pontryagin
  & $0.000$ & $0.000$ & $0.000$ & $0.000$
  & $0.000$ & $-97.7 \pm\;\;1.1$ \\
\bottomrule
\end{tabular}}
\end{table}

Several observations stand out.

\textbf{Otsu binarisation improves all IoU values.}
Compared to the previous fixed-threshold evaluation
(Euclidean GT IoU $= 0.287$, Hyperbolic $= 0.280$), the adaptive threshold
(Otsu $\approx 0.29$--$0.32$) yields substantially higher values
(Euclidean $0.658$, Hyperbolic $0.549$), confirming that fixed-threshold
comparisons across models with different activation scales are misleading.

\textbf{Euclidean achieves the best spatial alignment.}
The Euclidean head leads on GT IoU, Recall, and Pearson correlation, with
the Hyperbolic head showing higher GT Precision ($0.874$ vs $0.827$) —
suggesting that the hyperbolic embedding produces fewer false-positive
activations outside the sugarcane region, at the cost of reduced coverage.

\textbf{Pontryagin CAM degenerates on the first image.}
The mean activation of $0.000$ and all-zero GT metrics indicate that the
SegScoreCAM for the Pontryagin head produces a null heatmap (all feature
maps have zero or near-zero softmax contribution) on this image.  This is
consistent with the earlier finding of high mean activation ($0.366$)
across the dataset: the Pontryagin embedding yields \emph{bimodal}
heatmaps — either saturated (diffuse, high mean) or collapsed (null).

\textbf{Confidence gain is universally negative.}
All three models show a negative mean confidence change when the input is
masked by the sugarcane CAM.  This reveals that the models rely on
\emph{contextual information} beyond the highlighted region to maintain
their predictions.  The Hyperbolic model is least affected ($-40.6\%$),
suggesting its embedding captures more locally self-sufficient features.
The Pontryagin model shows the most severe drop ($-97.7\%$), consistent
with its collapsed heatmaps.  Importantly, the negative confidence gain is
\emph{not} a failure of the CAM method per se — it reflects a property
of the underlying model, namely that sugarcane predictions are
context-dependent at the scale of the $512 \times 512$ image.

% =============================================================================
\section{Discussion}
% =============================================================================

\paragraph{Decoupling $\kap$ and $\dpoly$.}
The separation of the Pontryagin signature $\kap$ from the SRF polynomial
degree $\dpoly$ is a key design correction.  Previously both were set to the
same integer, conflating three roles simultaneously: (i) the number of SRF
output components, (ii) the number of negative dimensions in $J$, and (iii)
the nonlinearity of $K_-$.  With independent parameters, the HPO can
identify whether a modest negative subspace ($\kap = 2$) paired with a
high-degree kernel ($\dpoly = 4$) outperforms a high-dimensional subspace
($\kap = 32$) with a linear kernel ($\dpoly = 1$) — something impossible to
test under the old parameterisation.

\paragraph{Expanded signature search space.}
The HPO range for $\kap$ has been extended from $\{1,2,3,4\}$ to
$\{1,\ldots,64\}$ with a log-uniform prior.  The previous restriction was
overly conservative: with $C_{\text{out}} = 64$ backbone channels and
$n_{\text{rff}} = m \cdot 64$, a signature of $\kap = 4$ produces a
Pontryagin space that is strongly dominated by the positive subspace.
Larger values of $\kap$ allow the indefinite structure to play a more
meaningful geometric role.  The log-uniform prior ensures that intermediate
values (e.g.\ $\kap = 8$ or $16$) are not underexplored relative to the
extremes.

\paragraph{Interpretability findings.}
Three concrete conclusions emerge from the expanded metric suite:

\begin{enumerate}[noitemsep]
  \item \textbf{Fixed-threshold CAM evaluation is unreliable.}
        The Otsu-adaptive threshold exposes substantially higher GT IoU
        ($0.658$ vs $0.287$ for Euclidean) compared to the fixed-$0.5$
        baseline, because models with small mean activation require a
        lower binarisation threshold to capture their concentrated responses.

  \item \textbf{Hyperbolic embeddings are more locally self-sufficient.}
        The Hyperbolic head achieves the highest GT Precision ($0.874$) and
        the smallest absolute confidence drop ($-40.6\%$), suggesting that
        the Poincaré-ball embedding concentrates discriminative information
        more tightly in the highlighted region than the Euclidean or
        Pontryagin heads.

  \item \textbf{Pontryagin CAMs collapse under the current HPO regime.}
        Null heatmaps (all-zero activation) on individual test images, combined
        with a high mean activation across the dataset, indicate bimodal
        behaviour: the Pontryagin head alternates between diffuse and
        collapsed activations depending on the image.  This instability is
        likely driven by the interaction of the cone, topographic, and balance
        penalties pushing embeddings away from the null cone in a way that
        disrupts spatial concentration.  The expanded $\kap$ search and the
        now-independent $\dpoly$ are expected to resolve this in subsequent
        HPO runs.
\end{enumerate}

% =============================================================================
\section{Conclusion}
% =============================================================================

We presented a semantic segmentation framework based on Pontryagin space
embeddings, combining Random Fourier Features (RFF) and Spherical Random
Features (SRF) with an indefinite J-metric and a J-hyperplane classifier.
Three heads — Euclidean, Hyperbolic, and Pontryagin — were trained on a UAV
sugarcane segmentation dataset under a controlled protocol (identical
backbone, optimiser, and schedule).

Key contributions are:
\begin{enumerate}[noitemsep]
  \item A pixel-level Pontryagin MLR head with composite loss (CE + cone +
        topographic + balance penalties).
  \item Decoupling of the Pontryagin signature $\kap$ from the SRF polynomial
        degree $\dpoly$, enabling independent HPO of both.
  \item SegScoreCAM: an adaptation of ScoreCAM for segmentation models,
        with an expanded metric suite comprising Otsu-adaptive binarisation,
        GT IoU, recall, precision, Pearson correlation, and pixel-level
        confidence gain on ground-truth class pixels.
  \item An interpretability comparison revealing that (a) fixed-threshold CAM
        evaluation substantially underestimates GT alignment, (b) the
        hyperbolic head achieves the highest GT precision ($0.874$) and
        smallest confidence drop, and (c) the Pontryagin head exhibits bimodal
        activation behaviour that is expected to stabilise after expanded HPO.
\end{enumerate}

Future work includes a full ablation over $(\kap, \dpoly)$ combinations,
analysis of the embedding geometry via t-SNE/UMAP in Pontryagin space,
investigation of the Pontryagin CAM instability under various composite-loss
weight schedules, and extension to multi-class segmentation.

% =============================================================================
\bibliographystyle{plain}
\begin{thebibliography}{9}

\bibitem{rahimi2007}
A. Rahimi and B. Recht.
\textit{Random Features for Large-Scale Kernel Machines}.
NeurIPS, 2007.

\bibitem{pennington2015}
J. Pennington, F. Yu, and S. Kumar.
\textit{Spherical Random Features for Polynomial Kernels}.
NeurIPS, 2015.

\bibitem{ganea2018}
O.-E. Ganea, G. Bécigneul, and T. Hofmann.
\textit{Hyperbolic Neural Networks}.
NeurIPS, 2018.

\bibitem{atigh2022}
M. G. Atigh, J. Schoep, E. Acar, N. van Noord, and P. Mettes.
\textit{Hyperbolic Image Segmentation}.
CVPR, 2022.

\bibitem{scorecam2020}
H. Wang, Z. Wang, M. Du, F. Yang, Z. Zhang, S. Ding, P. Mardziel, and X. Hu.
\textit{Score-CAM: Score-Weighted Visual Explanations for Convolutional Neural Networks}.
CVPR Workshops, 2020.

\end{thebibliography}

\end{document}
